\chapter{Glossary}
\label{chap:Glossary}

\textbf{Architectural layer:} A classification of firmware symbols into one of four categories --- application, middleware, low-level driver, or system runtime --- used to attribute memory costs to comparable functional blocks across vendors.\\

\textbf{Board Support Package (BSP):} A set of drivers and configuration files that adapt a vendor's generic HAL to a specific evaluation or development board.\\

\textbf{Context switch:} The process by which an RTOS saves the CPU state of the currently running task and restores the state of the next task to be executed. Context-switch efficiency directly affects multi-threaded throughput.\\

\textbf{Cross-vendor benchmark:} A controlled experiment in which the same workload runs on two vendor platforms under identical conditions (clock, optimisation level, scenario logic), producing comparable metrics.\\

\textbf{Firmware footprint:} The amount of non-volatile (flash/ROM) and volatile (RAM) memory a firmware image occupies once compiled and linked. It is a static measurement taken from the linked binary, not a dynamic runtime observation.\\

\textbf{Linker map:} A text file produced by the linker at the end of the build process. It lists every module and function linked into the final image together with its code size, constant-data size, and variable-data size. It is the primary input to the footprint analysis tool.\\

\textbf{Middleware:} Software libraries that sit between the hardware abstraction layer and the application, providing protocol or service logic (e.g.\ USB device stack, FreeRTOS kernel, FatFS, LwIP).\\

\textbf{Semantic equivalence:} The property that two benchmark implementations --- one on the reference platform and one on the competitor --- exhibit identical observable behaviour (task structure, timing, synchronisation, reporting format) so that their comparison is fair.\\

\textbf{Software stack:} A vertical slice of firmware functionality (e.g.\ USB device, FreeRTOS integration) comprising application code, middleware, and hardware drivers, benchmarked as a unit.\\

\textbf{Tick (RTOS):} The periodic timer interrupt that drives the RTOS scheduler. A 1\,ms tick means the scheduler evaluates task readiness every millisecond.\\

\pagebreak